\synctex=1
\documentclass[11pt,aspectratio=169]{beamer}

\usepackage[T1]{fontenc}
\usepackage[french,provide=*]{babel}
\usepackage{amsmath,amssymb,amsfonts}
\usepackage{mathtools}
\usepackage{bm}

% Police serif pour les mathématiques
\usefonttheme[onlymath]{serif}

% Git hash
\usepackage{xstring}
\usepackage{catchfile}
\immediate\write18{git rev-parse HEAD > git.hash}
\CatchFileDef{\HEAD}{git.hash}{\endlinechar=-1}
\newcommand{\gitrevision}{\StrLeft{\HEAD}{7}}

\usepackage{ccicons}

% Pied de page personnalisé
\setbeamertemplate{footline}{
  \hfill\vspace*{1pt}\href{https://creativecommons.org/publicdomain/zero/1.0/deed.fr}{\cczero}\enspace\href{https://github.com/stepan-a/time-series}{\gitrevision}\enspace--\enspace\insertframenumber/\inserttotalframenumber\enspace--\enspace\today\enspace}

% Supprimer la barre de navigation
\setbeamertemplate{navigation symbols}{}

% Commandes personnalisées
\newcommand{\E}{\mathbb{E}}
\newcommand{\Var}{\mathbb{V}}
\newcommand{\Cov}{\mathrm{Cov}}
\newcommand{\Corr}{\mathrm{Corr}}
\newcommand{\R}{\mathbb{R}}
\newcommand{\Z}{\mathbb{Z}}
\newcommand{\N}{\mathbb{N}}

\title{Processus Stochastiques Stationnaires}
\subtitle{Séries Temporelles}
\author{}
\date{}
\institute{}

\AtBeginSection[]
{
  \begin{frame}
    \frametitle{Plan}
    \tableofcontents[currentsection]
  \end{frame}
}

\begin{document}

\begin{frame}
\titlepage
\end{frame}

\begin{frame}{Plan}
\tableofcontents
\end{frame}

%==============================================================================
\section{Introduction aux processus stochastiques}
%==============================================================================

\begin{frame}{Définition d'un processus stochastique}

\begin{itemize}

\item Un \textbf{processus stochastique} est une suite de variables aléatoires réelles $\{X_t\}_{t \in \Z}$.\newline

\item Une \textbf{série temporelle} est une réalisation d'un processus stochastique.\newline

\item Pour caractériser complètement un processus stochastique, il faut spécifier la distribution jointe de $(X_{t_1}, X_{t_2}, \ldots, X_{t_n})$ pour tout $(t_1, t_2, \ldots, t_n) \in \Z^n$ et pour tout $n \in \N$, avec $t_i \neq t_j$ pour tout $i \neq j$.\newline

\item Il faut énormément d'information pour définir complètement un processus stochastique.

\end{itemize}

\end{frame}

\begin{frame}{Processus linéaires}

\begin{itemize}

\item \textbf{Problème de l'inférence :} Comment estimer les paramètres du processus stochastique à partir d'une seule réalisation (série temporelle) ?\newline

\item Pour traiter ce problème, on se restreint aux \textbf{processus linéaires} : processus dont les distributions jointes sont caractérisées par les moments d'ordre 1 et 2 :
  \[
  \mu_t = \E[X_t]
  \]
  \[
  \gamma(t,s) = \Cov(X_t, X_s) = \E\left[(X_t - \mu_t)(X_s - \mu_s)\right]
  \]

\end{itemize}

\end{frame}

\begin{frame}{Propriété de la fonction d'autocovariance}

\begin{itemize}

\item \textbf{Exercice :} Montrer que $\gamma(t,s) = \E[X_t X_s] - \mu_t \mu_s$.\newline

\item Démonstration :
  \begin{align*}
  \gamma(t,s) &= \E\left[(X_t - \mu_t)(X_s - \mu_s)\right] \\
  &= \E\left[X_t X_s - X_t \mu_s - \mu_t X_s + \mu_t \mu_s\right] \\
  &= \E[X_t X_s] - \E[X_t]\mu_s - \mu_t \E[X_s] + \mu_t \mu_s \\
  &= \E[X_t X_s] - \mu_t \mu_s - \mu_t \mu_s + \mu_t \mu_s \\
  &= \E[X_t X_s] - \mu_t \mu_s \quad \square
  \end{align*}

\end{itemize}

\end{frame}

\begin{frame}{Le problème de l'estimation avec une seule réalisation}

\begin{itemize}

\item \textbf{Exemple industriel :} Supposons qu'un processus stochastique décrit le diamètre d'un câble sortant d'une machine à intervalles réguliers (chaque mètre).\newline

\item Une série temporelle = mesures sur un câble de 1 km (1000 observations).\newline

\item Si l'on fabrique $K$ câbles avec la même machine, on peut estimer $\mu_t$ :
  \[
  \hat{\mu}_t = \frac{1}{K} \sum_{j=1}^{K} X_t^{(j)}
  \]
  où $X_t^{(j)}$ est le diamètre du câble $j$ au $t$-ème mètre.\newline

\item \textbf{Problème pour l'économiste :} On n'observe qu'une seule série temporelle pour le PIB... On ne peut pas « remonter le temps » pour obtenir une nouvelle réalisation.

\end{itemize}

\end{frame}

%==============================================================================
\section{Stationnarité}
%==============================================================================

\begin{frame}{Définition de la stationnarité (au sens strict)}

\begin{itemize}

\item Le processus stochastique $(X_t, t \in \Z)$ est \textbf{stationnaire} si et seulement si la distribution jointe de $(X_{t_1}, X_{t_2}, \ldots, X_{t_n})$ est identique à la distribution jointe de $(X_{t_1+h}, X_{t_2+h}, \ldots, X_{t_n+h})$ pour tout $(t_1, t_2, \ldots, t_n) \in \Z^n$, pour tout $n \in \N$, pour tout $h \in \Z$, avec $t_i \neq t_j$ pour $i \neq j$.\newline

\item La distribution jointe est \textbf{invariante dans le temps}.\newline

\item Pour les processus linéaires, nous n'avons pas besoin d'une définition si générale.

\end{itemize}

\end{frame}

\begin{frame}{Stationnarité au second ordre}

\begin{itemize}

\item Le processus stochastique $(X_t, t \in \Z)$ est dit \textbf{stationnaire au second ordre} si et seulement si ses moments d'ordre 1 et 2 sont invariants :
  \[
  \E[X_t] = \mu \quad \forall t
  \]
  \[
  \Var[X_t] = \gamma_0 \quad \forall t
  \]
  \[
  \Cov(X_t, X_s) = \gamma(t-s)
  \]\newline

\item Un processus stochastique $(X_t, t \in \Z)$ stationnaire au second ordre appartient à l'espace des variables aléatoires de carré intégrable $(L_2)$. La norme dans cet espace est définie par :
  \[
  \|X\| = \sqrt{\E[X^2]}
  \]

\end{itemize}

\end{frame}

\begin{frame}{Conséquences de la stationnarité au second ordre}

\begin{itemize}

\item Si le processus est stationnaire au second ordre, alors :
  \[
  \Cov(X_t, X_{t-h}) = \gamma(h) \quad \forall t
  \]
  La covariance ne dépend pas du temps mais seulement de la \textbf{distance} $h$ entre $X_t$ et $X_s$.\newline

\item En nous restreignant aux processus stochastiques stationnaires, nous avons considérablement diminué le nombre de paramètres à estimer.\newline

\item Mais pour pouvoir effectivement estimer les paramètres (par exemple l'espérance), il faut aussi que le processus soit \textbf{ergodique}.

\end{itemize}

\end{frame}

\begin{frame}{Ergodicité : intuition}

\begin{itemize}

\item Le concept d'ergodicité est techniquement difficile, mais l'idée intuitive est que des variables aléatoires doivent être d'autant moins corrélées qu'elles sont éloignées dans le temps.\newline

\item Si on utilise une moyenne empirique pour estimer l'espérance :
  \[
  \bar{X}_n = \frac{1}{n} \sum_{t=1}^{n} X_t
  \]
  alors $\E[\bar{X}_n] = \mu$ (sans biais) et surtout $\Var[\bar{X}_n] \xrightarrow[n \to \infty]{} 0$, assurant la convergence de $\bar{X}_n$ vers $\mu$.\newline

\item \textbf{Question clé :} Sous quelles conditions sur $\gamma(h)$ a-t-on $\Var[\bar{X}_n] \to 0$ ?

\end{itemize}

\end{frame}

\begin{frame}{Variance de la moyenne empirique (1/4)}

\begin{itemize}

\item Calculons la variance de $\bar{X}_n = \frac{1}{n} \sum_{t=1}^{n} X_t$ :
  \[
  \Var[\bar{X}_n] = \Var\left[\frac{1}{n} \sum_{t=1}^{n} X_t\right] = \frac{1}{n^2} \Var\left[\sum_{t=1}^{n} X_t\right]
  \]\newline

\item En développant la variance de la somme :
  \begin{align*}
  \Var\left[\sum_{t=1}^{n} X_t\right] &= \E\left[\left(\sum_{t=1}^{n} (X_t - \mu)\right)^2\right] = \sum_{t=1}^{n} \sum_{s=1}^{n} \Cov(X_t, X_s)
  \end{align*}\newline

\item Par stationnarité, $\Cov(X_t, X_s) = \gamma(t-s)$, donc :
  \[
  \Var\left[\sum_{t=1}^{n} X_t\right] = \sum_{t=1}^{n} \sum_{s=1}^{n} \gamma(t-s)
  \]

\end{itemize}

\end{frame}

\begin{frame}{Variance de la moyenne empirique (2/4)}

\begin{itemize}

\item On a $\displaystyle\Var\left[\sum_{t=1}^{n} X_t\right] = \sum_{t=1}^{n} \sum_{s=1}^{n} \gamma(t-s)$.\newline

\item Posons $h = t - s$. Pour chaque valeur de $h \in \{-(n-1), \ldots, n-1\}$, comptons le nombre de couples $(t,s)$ tels que $t - s = h$ :
  \begin{itemize}
    \item Si $h = 0$ : les couples $(1,1), (2,2), \ldots, (n,n)$ $\Rightarrow$ $n$ termes\newline
    \item Si $h = 1$ : les couples $(2,1), (3,2), \ldots, (n,n-1)$ $\Rightarrow$ $n-1$ termes\newline
    \item Plus généralement, pour $h \in \{-(n-1), \ldots, n-1\}$ : $n - |h|$ termes
  \end{itemize}

\item Donc :
  \[
  \sum_{t=1}^{n} \sum_{s=1}^{n} \gamma(t-s) = \sum_{h=-(n-1)}^{n-1} (n - |h|) \gamma(h)
  \]

\end{itemize}

\end{frame}

\begin{frame}{Variance de la moyenne empirique (3/4)}

\begin{itemize}

\item Nous avons donc :
  \[
  \Var[\bar{X}_n] = \frac{1}{n^2} \sum_{h=-(n-1)}^{n-1} (n - |h|) \gamma(h)
  \]\newline

\item En utilisant la parité $\gamma(h) = \gamma(-h)$ :
  \begin{align*}
  \Var[\bar{X}_n] &= \frac{1}{n^2} \left[ n\gamma(0) + 2\sum_{h=1}^{n-1} (n-h) \gamma(h) \right]
  \end{align*}\newline

\item Résultat :
  \[
  \boxed{\Var[\bar{X}_n] = \frac{\gamma(0)}{n} + \frac{2}{n} \sum_{h=1}^{n-1} \left(1 - \frac{h}{n}\right) \gamma(h)}
  \]

\end{itemize}

\end{frame}

\begin{frame}{Variance de la moyenne empirique (4/4)}

\begin{itemize}

\item \textbf{Cas 1 (bruit blanc) :} Si $\gamma(h) = 0$ pour $h \neq 0$ :
  \[
  \Var[\bar{X}_n] = \frac{\gamma(0)}{n} = \frac{\sigma^2}{n} \xrightarrow[n \to \infty]{} 0 \quad \checkmark
  \]\newline

\item \textbf{Cas 2 (autocovariance constante) :} Si $\gamma(h) = c > 0$ pour tout $h$ :
  \begin{align*}
  \Var[\bar{X}_n] &= \frac{c}{n} + \frac{2c}{n} \left[ n - 1 - \frac{n-1}{2} \right] \xrightarrow[n \to \infty]{} c \neq 0 \quad \times
  \end{align*}\newline

\item La variance ne converge pas vers 0 : l'estimation est impossible.

\end{itemize}

\end{frame}

\begin{frame}{Condition suffisante de convergence}

\begin{itemize}

\item Si la fonction d'autocovariance est \textbf{absolument sommable} :
  \[
  \sum_{h=-\infty}^{+\infty} |\gamma(h)| < +\infty
  \]
  alors $\Var[\bar{X}_n] \to 0$ quand $n \to \infty$.\newline

\item Démonstration : Si $\sum_{h} |\gamma(h)| < \infty$, alors :
  \[
  \Var[\bar{X}_n] \leq \frac{1}{n} \sum_{h=-(n-1)}^{n-1} |\gamma(h)| \leq \frac{1}{n} \sum_{h=-\infty}^{+\infty} |\gamma(h)| \xrightarrow[n \to \infty]{} 0
  \]\newline

\item Cette condition implique que $\gamma(h) \to 0$ suffisamment vite quand $|h| \to \infty$. Elle est satisfaite par la plupart des modèles ARMA.

\end{itemize}

\end{frame}

\begin{frame}{Comportement asymptotique de la variance}

\begin{itemize}

\item Quand $n \to \infty$, on peut montrer que :
  \[
  n \cdot \Var[\bar{X}_n] \xrightarrow[n \to \infty]{} \sum_{h=-\infty}^{+\infty} \gamma(h) = 2\pi f(0)
  \]
  où $f(\omega)$ est la densité spectrale du processus.\newline

\item La quantité $\sum_{h} \gamma(h) = 2\pi f(0)$ est appelée \textbf{variance de long terme}.\newline

\item Interprétation :
  \begin{itemize}
    \item Si $\sum_{h} \gamma(h) > \gamma(0)$ : la dépendance positive augmente la variance de $\bar{X}_n$.\newline
    \item Si $\sum_{h} \gamma(h) < \gamma(0)$ : la dépendance négative diminue la variance de $\bar{X}_n$.
  \end{itemize}

\end{itemize}

\end{frame}

%==============================================================================
\section{Moments d'un processus stationnaire}
%==============================================================================

\begin{frame}{Fonction d'autocovariance}

\begin{itemize}

\item On note $\gamma : \Z \to \R$, $h \mapsto \gamma(h)$ la \textbf{fonction d'autocovariance} d'un processus stochastique stationnaire au second ordre.\newline

\item Propriétés :
  \begin{enumerate}
    \item $\gamma(0) = \Var[X_t] \geq 0$\newline
    \item $\gamma(h) = \gamma(-h)$ \quad (parité)\newline
    \item $|\gamma(h)| \leq \gamma(0)$ \quad pour tout $h$\newline
    \item $\gamma$ est une fonction \textbf{positive} (au sens de la positivité définie)
  \end{enumerate}

\end{itemize}

\end{frame}

\begin{frame}{Démonstration de la parité}

\begin{itemize}

\item \textbf{Propriété :} $\gamma(h) = \gamma(-h)$.\newline

\item Démonstration : Il suffit de remarquer que
  \begin{align*}
  \gamma(-h) &= \Cov(X_t, X_{t+h}) \\
  &= \E\left[(X_t - \mu)(X_{t+h} - \mu)\right] \\
  &= \E\left[(X_{t+h} - \mu)(X_t - \mu)\right] \\
  &= \gamma(h) \quad \square
  \end{align*}

\end{itemize}

\end{frame}

\begin{frame}{Inégalité de Cauchy-Schwarz pour la covariance (1/2)}

\begin{itemize}

\item Pour deux variables aléatoires $U$ et $V$ de carré intégrable :
  \[
  |\Cov(U, V)| \leq \sqrt{\Var(U)} \cdot \sqrt{\Var(V)}
  \]\newline

\item Démonstration : Sans perte de généralité, supposons $\E[U] = \E[V] = 0$. Alors $\Cov(U,V) = \E[UV]$ et $\Var(U) = \E[U^2]$.\newline

\item Pour tout $\lambda \in \R$, la variable aléatoire $(U + \lambda V)^2$ est positive, donc :
  \[
  \E\left[(U + \lambda V)^2\right] \geq 0 \quad \forall \lambda \in \R
  \]\newline

\item En développant :
  \[
  \E[U^2] + 2\lambda \E[UV] + \lambda^2 \E[V^2] \geq 0
  \]

\end{itemize}

\end{frame}

\begin{frame}{Inégalité de Cauchy-Schwarz pour la covariance (2/2)}

\begin{itemize}

\item Nous avons un trinôme en $\lambda$ :
  \[
  P(\lambda) = \E[V^2] \lambda^2 + 2\E[UV] \lambda + \E[U^2] \geq 0 \quad \forall \lambda \in \R
  \]\newline

\item Pour qu'un trinôme $a\lambda^2 + b\lambda + c$ soit toujours $\geq 0$, son discriminant doit être $\leq 0$ :
  \[
  \Delta = 4\E[UV]^2 - 4\E[U^2]\E[V^2] \leq 0
  \]\newline

\item En prenant la racine carrée et en revenant aux variables non centrées :
  \[
  \boxed{|\Cov(U, V)| \leq \sqrt{\Var(U)} \cdot \sqrt{\Var(V)}} \quad \square
  \]

\end{itemize}

\end{frame}

\begin{frame}{Démonstration de $|\gamma(h)| \leq \gamma(0)$}

\begin{itemize}

\item \textbf{Propriété :} $|\gamma(h)| \leq \gamma(0)$ pour tout $h$.\newline

\item Cette propriété découle de l'inégalité de Cauchy-Schwarz.\newline

\item En appliquant cette inégalité à $U = X_t$ et $V = X_{t-h}$ :
  \[
  |\gamma(h)| = |\Cov(X_t, X_{t-h})| \leq \sqrt{\Var(X_t)} \cdot \sqrt{\Var(X_{t-h})}
  \]\newline

\item Par stationnarité, $\Var(X_t) = \Var(X_{t-h}) = \gamma(0)$, donc :
  \[
  |\gamma(h)| \leq \sqrt{\gamma(0)} \cdot \sqrt{\gamma(0)} = \gamma(0) \quad \square
  \]

\end{itemize}

\end{frame}

\begin{frame}{Interprétation et conséquences}

\begin{itemize}

\item L'inégalité $|\gamma(h)| \leq \gamma(0)$ signifie que la covariance entre $X_t$ et $X_{t-h}$ ne peut jamais dépasser (en valeur absolue) la variance du processus.\newline

\item \textbf{Cas d'égalité :} $|\gamma(h)| = \gamma(0)$ si et seulement si $X_t$ et $X_{t-h}$ sont parfaitement corrélés.\newline

\item En divisant par $\gamma(0) > 0$, on obtient :
  \[
  \left|\frac{\gamma(h)}{\gamma(0)}\right| \leq 1 \quad \forall h
  \]\newline

\item Le ratio $\gamma(h)/\gamma(0)$ est toujours compris entre $-1$ et $1$. Ce ratio définit la \textbf{fonction d'autocorrélation} $\rho(h)$.

\end{itemize}

\end{frame}

\begin{frame}{Positivité de la fonction d'autocovariance}

\begin{itemize}

\item Dire que la fonction d'autocovariance est \textbf{positive} ne veut pas dire que $\gamma(h) \geq 0$ pour tout $h$. Cette propriété découle de la positivité de la variance.\newline

\item On sait que $\Var\left[\sum_{i=1}^n a_i X_{t_i}\right] \geq 0$ pour tout vecteur $(a_1, \ldots, a_n)$.\newline

\item Supposons $\E[X_{t_i}] = 0$ pour tout $i$. Alors :
  \begin{align*}
  \Var\left[\sum_{i=1}^n a_i X_{t_i}\right] &= \sum_{i=1}^n \sum_{j=1}^n a_i a_j \gamma(t_i - t_j)
  \end{align*}\newline

\item Cette double somme doit être $\geq 0$ pour tout vecteur $\mathbf{a} \neq \mathbf{0}$.

\end{itemize}

\end{frame}

%==============================================================================
\section{Construction de processus stationnaires}
%==============================================================================

\begin{frame}{Théorème de construction}

\begin{itemize}

\item Si $(X_t, t \in \Z)$ est un processus stochastique stationnaire et si $(a_i)_{i \geq 0}$ est une suite de nombres réels \textbf{absolument sommable} :
  \[
  \sum_{i=0}^{\infty} |a_i| < +\infty
  \]
  alors
  \[
  Y_t = \sum_{i=0}^{\infty} a_i X_{t-i}
  \]
  définit un nouveau processus stochastique stationnaire.

\end{itemize}

\end{frame}

\begin{frame}{Éléments de preuve (1/2)}

\begin{itemize}

\item Si $(X_t, t \in \Z)$ est de carré intégrable, alors il en va de même pour $(Y_t, t \in \Z)$.\newline

\item En effet :
  \[
  \|Y_t\| = \left\|\sum_{i=0}^{\infty} a_i X_{t-i}\right\| \leq \sum_{i=0}^{\infty} |a_i| \cdot \|X_{t-i}\| = \|X\| \sum_{i=0}^{\infty} |a_i| < +\infty
  \]\newline

\item L'espérance de $(Y_t, t \in \Z)$ est :
  \[
  \E[Y_t] = \E\left[\sum_{i=0}^{\infty} a_i X_{t-i}\right] = \sum_{i=0}^{\infty} a_i \E[X_{t-i}] = \mu \sum_{i=0}^{\infty} a_i = \mu_Y
  \]\newline

\item L'espérance est constante et ne dépend pas de $t$.

\end{itemize}

\end{frame}

\begin{frame}{Éléments de preuve (2/2)}

\begin{itemize}

\item La fonction d'autocovariance de $(Y_t, t \in \Z)$ est :
  \begin{align*}
  \gamma_Y(h) &= \Cov(Y_t, Y_{t-h}) \\
  &= \Cov\left(\sum_{i=0}^{\infty} a_i X_{t-i}, \sum_{j=0}^{\infty} a_j X_{t-h-j}\right) \\
  &= \sum_{i=0}^{\infty} \sum_{j=0}^{\infty} a_i a_j \gamma_X(h+i-j)
  \end{align*}\newline

\item Les moments d'ordre 1 et 2 sont bien indépendants du temps. \hfill $\square$

\end{itemize}

\end{frame}

%==============================================================================
\section{Exemple : Processus MA(1)}
%==============================================================================

\begin{frame}{Bruit blanc}

\begin{itemize}

\item Soit $(X_t, t \in \Z)$ une suite de variables aléatoires \textbf{indépendamment et identiquement distribuées} (bruit blanc) avec :
  \[
  \E[X_t] = 0 \quad \forall t
  \]
  \[
  \Var[X_t] = \sigma^2 \quad \forall t
  \]
  \[
  \gamma(h) = 0 \quad \forall h \neq 0
  \]\newline

\item Ce processus stochastique est stationnaire.\newline

\item On note souvent : $X_t \sim \text{BB}(0, \sigma^2)$ ou $\varepsilon_t \sim \text{BB}(0, \sigma^2)$.

\end{itemize}

\end{frame}

\begin{frame}{Définition du processus MA(1)}

\begin{itemize}

\item Définissons un nouveau processus stochastique $(Y_t, t \in \Z)$ :
  \[
  Y_t = \lambda X_t + \rho X_{t-1}
  \]\newline

\item En termes du théorème précédent : $a_0 = \lambda$, $a_1 = \rho$, $a_i = 0$ pour $i \notin \{0, 1\}$.\newline

\item On montre facilement qu'il s'agit d'un processus stochastique stationnaire au second ordre.

\end{itemize}

\end{frame}

\begin{frame}{Espérance du processus MA(1)}

\begin{itemize}

\item L'espérance est nulle pour tout $t$ :
  \begin{align*}
  \E[Y_t] &= \E[\lambda X_t + \rho X_{t-1}] \\
  &= \lambda \E[X_t] + \rho \E[X_{t-1}] \\
  &= \lambda \cdot 0 + \rho \cdot 0 \\
  &= 0
  \end{align*}

\end{itemize}

\end{frame}

\begin{frame}{Variance du processus MA(1)}

\begin{itemize}

\item La variance est finie et constante :
  \begin{align*}
  \Var[Y_t] &= \Var[\lambda X_t + \rho X_{t-1}] \\
  &= \E\left[(\lambda X_t + \rho X_{t-1})^2\right] \quad \text{(car } \E[Y_t] = 0\text{)} \\
  &= \E\left[\lambda^2 X_t^2 + 2\lambda\rho X_t X_{t-1} + \rho^2 X_{t-1}^2\right] \\
  &= \lambda^2 \E[X_t^2] + 2\lambda\rho \underbrace{\E[X_t X_{t-1}]}_{= 0 \text{ (indép.)}} + \rho^2 \E[X_{t-1}^2] \\
  &= (\lambda^2 + \rho^2) \sigma^2
  \end{align*}

\end{itemize}

\end{frame}

\begin{frame}{Fonction d'autocovariance du MA(1)}

\begin{itemize}

\item Calculons la fonction d'autocovariance $\gamma(h) = \Cov(Y_t, Y_{t-h})$ :
  \begin{align*}
  \gamma(h) &= \Cov(\lambda X_t + \rho X_{t-1}, \lambda X_{t-h} + \rho X_{t-h-1}) \\
  &= \lambda^2 \Cov(X_t, X_{t-h}) + \lambda\rho \Cov(X_t, X_{t-h-1}) \\
  &\quad + \rho\lambda \Cov(X_{t-1}, X_{t-h}) + \rho^2 \Cov(X_{t-1}, X_{t-h-1})
  \end{align*}\newline

\item Au total :
  \[
  \gamma(h) = \begin{cases}
  (\lambda^2 + \rho^2)\sigma^2 & \text{si } h = 0 \\
  \lambda\rho\sigma^2 & \text{si } h = \pm 1 \\
  0 & \text{sinon}
  \end{cases}
  \]

\end{itemize}

\end{frame}

\begin{frame}{Interprétation}

\begin{itemize}

\item Contrairement au processus d'origine $(X_t, t \in \Z)$ (bruit blanc), le processus $(Y_t, t \in \Z)$ admet de la \textbf{dépendance} (relativement limitée).\newline

\item La covariance entre $Y_t$ et $Y_{t-1}$ est non nulle : $\gamma(1) = \lambda\rho\sigma^2$.\newline

\item Mais cette dépendance est de \textbf{courte mémoire} : $\gamma(h) = 0$ pour $|h| > 1$.\newline

\item Ce processus est appelé \textbf{MA(1)} (moyenne mobile d'ordre 1).

\end{itemize}

\end{frame}

%==============================================================================
\section{Fonction d'autocorrélation}
%==============================================================================

\begin{frame}{Innovation d'un processus}

\begin{itemize}

\item Si $(X_t, t \in \Z)$ est un processus stochastique stationnaire au second ordre, on définit son \textbf{innovation} par :
  \[
  \varepsilon_t = X_t - X_t^*
  \]
  où $X_t^*$ est la régression linéaire de $X_t$ sur son passé $(X_s, s < t)$.\newline

\item $X_t^*$ est la meilleure prévision (linéaire) de $X_t$ basée sur l'ensemble d'information $I_t = \{X_{t-1}, X_{t-2}, \ldots\}$.\newline

\item Par construction, l'innovation $\varepsilon_t$ est non corrélée avec le passé de $X$ : elle peut être interprétée comme une \textbf{erreur de prévision}.\newline

\item On peut montrer que si $(X_t, t \in \Z)$ est stationnaire au second ordre, alors son innovation est un \textbf{bruit blanc}.

\end{itemize}

\end{frame}

\begin{frame}{Fonction d'autocorrélation}

\begin{itemize}

\item La \textbf{fonction d'autocorrélation} est définie par :
  \[
  \rho(h) = \frac{\gamma(h)}{\gamma(0)}
  \]\newline

\item $\rho(h)$ mesure la corrélation entre $X_t$ et $X_{t+h}$ (ou $X_{t-h}$) :
  \[
  \rho(h) = \Corr(X_t, X_{t+h})
  \]\newline

\item La fonction d'autocorrélation est une normalisation de la fonction d'autocovariance. Elle hérite de ses propriétés : parité et positivité.

\end{itemize}

\end{frame}

\begin{frame}{Exemple}

\begin{itemize}

\item Soit $(\varepsilon_t, t \in \Z)$ un bruit blanc. On définit le processus $(X_t, t \in \Z)$ :
  \[
  X_t = \varepsilon_t - \varepsilon_{t-12}
  \]\newline

\item La fonction d'autocovariance est :
  \[
  \gamma(h) = \begin{cases}
  2\sigma^2 & \text{si } h = 0 \\
  -\sigma^2 & \text{si } h = \pm 12 \\
  0 & \text{sinon}
  \end{cases}
  \]\newline

\item La fonction d'autocorrélation est :
  \[
  \rho(h) = \begin{cases}
  1 & \text{si } h = 0 \\
  -\frac{1}{2} & \text{si } h = \pm 12 \\
  0 & \text{sinon}
  \end{cases}
  \]

\end{itemize}

\end{frame}

%==============================================================================
\section{Matrice d'autocorrélation et contraintes}
%==============================================================================

\begin{frame}{Matrice d'autocorrélation}

\begin{itemize}

\item La \textbf{matrice d'autocorrélation} d'ordre $m$ contient les corrélations entre $m$ $X$ successifs : $X_t, X_{t+1}, \ldots, X_{t+m-1}$.
  \[
  R(m) = \begin{pmatrix}
  1 & \rho(1) & \rho(2) & \cdots & \rho(m-1) \\
  \rho(1) & 1 & \rho(1) & \cdots & \rho(m-2) \\
  \rho(2) & \rho(1) & 1 & \cdots & \rho(m-3) \\
  \vdots & \vdots & \vdots & \ddots & \vdots \\
  \rho(m-1) & \rho(m-2) & \rho(m-3) & \cdots & 1
  \end{pmatrix}
  \]\newline

\item La positivité de la fonction $\rho(h)$ implique que $R(m)$ est une \textbf{matrice définie positive} :
  \[
  \mathbf{a}' R(m) \mathbf{a} > 0 \quad \forall \mathbf{a} \in \R^m, \mathbf{a} \neq \mathbf{0}
  \]

\end{itemize}

\end{frame}

\begin{frame}{Contraintes sur l'autocorrélation d'ordre 1}

\begin{itemize}

\item $|R(m)| > 0$ pour tout $m \in \N^*$.\newline

\item Application pour $m = 2$ :
  \[
  |R(2)| > 0 \Leftrightarrow \begin{vmatrix} 1 & \rho(1) \\ \rho(1) & 1 \end{vmatrix} > 0 \Leftrightarrow 1 - \rho(1)^2 > 0
  \]\newline

\item L'autocorrélation d'ordre 1 doit être strictement inférieure à 1 en valeur absolue :
  \[
  |\rho(1)| < 1
  \]

\end{itemize}

\end{frame}

\begin{frame}{Contraintes sur l'autocorrélation d'ordre 2}

\begin{itemize}

\item Application pour $m = 3$ :
  \[
  |R(3)| > 0 \Leftrightarrow \begin{vmatrix} 1 & \rho(1) & \rho(2) \\ \rho(1) & 1 & \rho(1) \\ \rho(2) & \rho(1) & 1 \end{vmatrix} > 0
  \]\newline

\item Après calcul du déterminant, on obtient la contrainte :\newline

\item Les valeurs possibles de $\rho(2)$ sont contraintes par $\rho(1)$ :
  \[
  \rho(2) \in \left[2\rho(1)^2 - 1, 1\right)
  \]

\end{itemize}

\end{frame}

%==============================================================================
\section{Autocorrélation partielle}
%==============================================================================

\begin{frame}{Prévision linéaire}

\begin{itemize}

\item Soit $(X_t, t \in \Z)$ un processus stochastique stationnaire au second ordre.\newline

\item On s'intéresse à la meilleure prévision linéaire de $X_t$ étant données les $K$ valeurs précédentes : $X_{t-1}, X_{t-2}, \ldots, X_{t-K}$.\newline

\item En supposant $\E[X_t] = 0$ :
  \[
  \E[X_t | X_{t-1}, \ldots, X_{t-K}] = a_1^{(K)} X_{t-1} + a_2^{(K)} X_{t-2} + \cdots + a_K^{(K)} X_{t-K}
  \]\newline

\item Le vecteur $\mathbf{a}^{(K)} = (a_1^{(K)}, a_2^{(K)}, \ldots, a_K^{(K)})'$ est donné par :
  \[
  \mathbf{a}^{(K)} = R(K)^{-1} \begin{pmatrix} \rho(1) \\ \rho(2) \\ \vdots \\ \rho(K) \end{pmatrix}
  \]

\end{itemize}

\end{frame}

\begin{frame}{Coefficient de corrélation partielle}

\begin{itemize}

\item Le \textbf{coefficient de corrélation partielle} est :
  \[
  r(K) = a_K^{(K)}
  \]
  Il mesure le lien entre $X_t$ et $X_{t-K}$ une fois que l'on a purgé l'effet de $X_{t-1}, X_{t-2}, \ldots, X_{t-K+1}$.\newline

\item Les fonctions d'autocorrélation et d'autocorrélation partielle sont \textbf{équivalentes} :
  \begin{itemize}
    \item $r(K)$ est une fonction de $(\rho(1), \ldots, \rho(K))$\newline
    \item Inversement, on peut déduire $\rho(K)$ de $(r(1), r(2), \ldots, r(K))$
  \end{itemize}

\item Pour $K = 1$ : $a_1^{(1)} = r(1) = \rho(1)$.

\end{itemize}

\end{frame}

\begin{frame}{Expression explicite de l'autocorrélation partielle}

\begin{itemize}

\item En utilisant la structure par blocs de la matrice, on peut montrer :
  \[
  r(1) = \rho(1)
  \]
  \[
  r(2) = \frac{\rho(2) - \rho(1)^2}{1 - \rho(1)^2}
  \]\newline

\item Plus généralement, on peut obtenir une expression explicite de $r(K)$ en exploitant la structure de la matrice d'autocorrélation.

\end{itemize}

\end{frame}

%==============================================================================
\section{Densité spectrale}
%==============================================================================

\begin{frame}{Introduction à l'approche fréquentielle}

\begin{itemize}

\item Une façon alternative de caractériser un processus stochastique est d'étudier la \textbf{densité spectrale}.\newline

\item Cette approche est équivalente à la fonction d'autocovariance.\newline

\item Plutôt que d'aborder un processus stochastique dans le \textbf{domaine temporel}, on s'intéresse au \textbf{domaine des fréquences}.\newline

\item On considère un processus de la forme :
  \[
  X_t = \sum_{i=0}^{\infty} a_i \varepsilon_{t-i}
  \]
  où $(\varepsilon_t, t \in \Z) \sim \text{BB}(0, \sigma^2)$ et $(a_i)$ est absolument sommable.

\end{itemize}

\end{frame}

\begin{frame}{Fonction d'autocovariance du processus $X$}
\textbf{Calcul de $\gamma(h)$ :} Pour $h \geq 0$,
\begin{align*}
\gamma(h) &= \Cov(X_t, X_{t-h}) = \Cov\left( \sum_{i=0}^{\infty} a_i \varepsilon_{t-i}, \sum_{j=0}^{\infty} a_j \varepsilon_{t-h-j} \right) \\
&= \sum_{i=0}^{\infty} \sum_{j=0}^{\infty} a_i a_j \Cov(\varepsilon_{t-i}, \varepsilon_{t-h-j})
\end{align*}

La covariance $\Cov(\varepsilon_{t-i}, \varepsilon_{t-h-j})$ est non nulle uniquement si $t-i = t-h-j$, i.e., $j = i - h$.

\vspace{0.3cm}

Donc pour $h \geq 0$ :
\[
\gamma(h) = \sigma^2 \sum_{i=h}^{\infty} a_i a_{i-h} = \sigma^2 \sum_{i=0}^{\infty} a_{i+h} a_i
\]

Par parité : $\gamma(-h) = \gamma(h)$.
\end{frame}

\begin{frame}{Sommabilité absolue de $\gamma(h)$}
\begin{block}{Théorème}
Si $(a_i)$ est absolument sommable, alors $\gamma(h)$ est absolument sommable.
\end{block}

\textbf{Démonstration :}
\begin{align*}
\sum_{h=-\infty}^{+\infty} |\gamma(h)| &= \sum_{h=-\infty}^{+\infty} \left| \sigma^2 \sum_{i=0}^{\infty} a_i a_{i+|h|} \right| \\
&\leq \sigma^2 \sum_{h=-\infty}^{+\infty} \sum_{i=0}^{\infty} |a_i| |a_{i+|h|}|
\end{align*}

En échangeant les sommes et en posant $j = i + |h|$ :
\[
\sum_{h=-\infty}^{+\infty} |\gamma(h)| \leq \sigma^2 \sum_{i=0}^{\infty} |a_i| \sum_{j=0}^{\infty} |a_j| = \sigma^2 \left( \sum_{i=0}^{\infty} |a_i| \right)^2 < +\infty
\]

La fonction d'autocovariance doit donc converger suffisament rapidement vers zéro.
\end{frame}

\begin{frame}{Définition de la densité spectrale}

\begin{itemize}

\item La \textbf{densité spectrale} du processus $(X_t, t \in \Z)$ est la fonction réelle définie par :
  \[
  f(\omega) = \frac{1}{2\pi} \sum_{h=-\infty}^{+\infty} \gamma(h) e^{-i\omega h}
  \]
  pour tout $\omega \in \R$.\newline

\item Cette fonction existe car $\gamma(h)$ est absolument sommable.\newline

\item On peut montrer que :
  \[
  f(\omega) = \frac{\gamma(0)}{2\pi} + \frac{1}{\pi} \sum_{h=1}^{\infty} \gamma(h) \cos(\omega h)
  \]\newline

\item La densité spectrale est une fonction continue, périodique et symétrique.

\end{itemize}

\end{frame}

\begin{frame}{La densité spectrale est à valeurs réelles}

\begin{itemize}

\item La définition $f(\omega) = \frac{1}{2\pi} \sum_{h=-\infty}^{+\infty} \gamma(h) e^{-i\omega h}$ fait intervenir des exponentielles complexes. Montrons que $f(\omega) \in \R$.\newline

\item On regroupe les termes symétriques et on utilise $\gamma(-h) = \gamma(h)$ :
  \[
  f(\omega) = \frac{1}{2\pi} \left[ \gamma(0) + \sum_{h=1}^{\infty} \gamma(h) \left( e^{-i\omega h} + e^{i\omega h} \right) \right]
  \]\newline

\item Or $e^{-i\omega h} + e^{i\omega h} = 2\cos(\omega h)$, d'où :
  \[
  \boxed{f(\omega) = \frac{\gamma(0)}{2\pi} + \frac{1}{\pi} \sum_{h=1}^{\infty} \gamma(h) \cos(\omega h) \in \R}
  \]

\end{itemize}

\end{frame}

\begin{frame}{Symétrie de la densité spectrale}

\begin{itemize}

\item La représentation en cosinus montre immédiatement que $f$ est une fonction \textbf{paire} :
  \[
  f(-\omega) = \frac{\gamma(0)}{2\pi} + \frac{1}{\pi} \sum_{h=1}^{\infty} \gamma(h) \cos(-\omega h) = f(\omega)
  \]
  car $\cos(-x) = \cos(x)$.\newline

\item De plus, $f$ est \textbf{$2\pi$-périodique} car $\cos(\omega h)$ est $2\pi$-périodique en $\omega$.\newline

\item En combinant parité et périodicité, il suffit de connaître $f$ sur $[0, \pi]$ pour la connaître partout.\newline

\item C'est pourquoi on représente généralement la densité spectrale sur l'intervalle $[0, \pi]$.

\end{itemize}

\end{frame}


\begin{frame}{Équivalence avec la fonction d'autocovariance}

\begin{itemize}

\item La densité spectrale et la fonction d'autocovariance sont \textbf{équivalentes}.\newline

\item De l'autocovariance à la densité spectrale : définition de $f(\omega)$.\newline

\item De la densité spectrale à l'autocovariance :
  \[
  \gamma(h) = \int_{-\pi}^{\pi} f(\omega) \cos(\omega h) \, d\omega
  \]\newline

\item En particulier, la variance s'exprime comme une intégrale de la densité spectrale :
  \[
  \gamma(0) = \int_{-\pi}^{\pi} f(\omega) \, d\omega
  \]

\end{itemize}

\end{frame}

\begin{frame}{Exemple : Densité spectrale d'un bruit blanc}

\begin{itemize}

\item Soit $(X_t, t \in \Z)$ une suite de variables aléatoires i.i.d. d'espérance nulle et de variance $\sigma^2$ (bruit blanc).\newline

\item La fonction d'autocovariance est :
  \[
  \gamma(h) = \begin{cases}
  \sigma^2 & \text{si } h = 0 \\
  0 & \text{sinon}
  \end{cases}
  \]\newline

\item La densité spectrale est donc :
  \[
  f(\omega) = \frac{\sigma^2}{2\pi}
  \]\newline

\item La densité spectrale est \textbf{constante} : toutes les fréquences contribuent également à la variance du processus. C'est pourquoi on parle de \textbf{bruit blanc}.

\end{itemize}

\end{frame}

\begin{frame}{Réciproque}

\begin{itemize}

\item Si la densité spectrale est plate (constante), alors le processus associé est un bruit blanc.\newline

\item Démonstration : Supposons que $f(\omega) = \kappa > 0$ pour tout $\omega$. Alors :
  \[
  \gamma(h) = \int_{-\pi}^{\pi} \kappa \cos(\omega h) \, d\omega = \begin{cases}
  \kappa \cdot 2\pi & \text{si } h = 0 \\
  0 & \text{sinon}
  \end{cases}
  \]\newline

\item $(X_t, t \in \Z)$ est bien un bruit blanc de variance $2\pi\kappa$. \hfill $\square$

\end{itemize}

\end{frame}

\begin{frame}{Interprétation de la densité spectrale}

\begin{itemize}

\item La densité spectrale $f(\omega)$ mesure la \textbf{contribution de la fréquence $\omega$} à la variance totale du processus :
  \[
  \gamma(0) = \int_{-\pi}^{\pi} f(\omega) \, d\omega = 2 \int_{0}^{\pi} f(\omega) \, d\omega
  \]\newline

\item Un \textbf{pic} dans $f(\omega)$ à la fréquence $\omega_0$ indique une composante cyclique dominante de période :
  \[
  T = \frac{2\pi}{\omega_0}\quad\text{(en nombre de périodes d'échantillonnage)}
  \]

\item Exemples :
  \begin{itemize}
    \item $\omega_0 = \pi$ $\Rightarrow$ $T = 2$ : oscillation à la fréquence maximale (alternance)\newline
    \item $\omega_0 = \pi/6$ $\Rightarrow$ $T = 12$ : cycle annuel pour des données mensuelles\newline
    \item $\omega_0 \approx 0$ $\Rightarrow$ $T \to \infty$ : composante de très basse fréquence (tendance)
  \end{itemize}

\end{itemize}

\end{frame}

\begin{frame}{Identification des cycles}

\begin{itemize}

\item En pratique, on estime $f(\omega)$ à partir des données (périodogramme) et on recherche les pics.\newline

\item Les \textbf{hautes fréquences} ($\omega$ proche de $\pi$) correspondent aux fluctuations rapides, de courte période.\newline

\item Les \textbf{basses fréquences} ($\omega$ proche de $0$) correspondent aux mouvements lents, de longue période.\newline

\item Un processus avec $f(\omega)$ concentré près de $\omega = 0$ est \textbf{persistant} : les chocs ont des effets durables.\newline

\item Un processus avec $f(\omega)$ concentré près de $\omega = \pi$ est \textbf{antipersistant} : tendance au retour à la moyenne rapide.

\end{itemize}

\end{frame}


\begin{frame}{Densité spectrale d'une transformation linéaire}

\begin{itemize}

\item Soit $(X_t, t \in \Z)$ un processus stationnaire de densité spectrale $f_X(\omega)$. Soit $(a_i)_{i \geq 0}$ une suite absolument sommable. Alors le processus $(Y_t, t \in \Z)$ défini par :
  \[
  Y_t = \sum_{i=0}^{\infty} a_i X_{t-i}
  \]
  a pour densité spectrale :
  \[
  f_Y(\omega) = \left|\sum_{i=0}^{\infty} a_i e^{-i\omega i}\right|^2 f_X(\omega)
  \]

\end{itemize}

\end{frame}

%==============================================================================
\section{Opérateur retard}
%==============================================================================

\begin{frame}{Définition de l'opérateur retard}

\begin{itemize}

\item L'\textbf{opérateur retard} $L$ transforme un processus $(X_t, t \in \Z)$ en un processus $(Y_t, t \in \Z)$ tel que :
  \[
  Y_t = LX_t = X_{t-1}
  \]\newline

\item Cet opérateur est linéaire et inversible.\newline

\item Son inverse est l'\textbf{opérateur avance} $F$ :
  \[
  Z_t = FX_t = X_{t+1}
  \]\newline

\item Par construction : $L \cdot F = F \cdot L = 1$.\newline

\item En appliquant plusieurs fois l'opérateur retard :
  \[
  L^k X_t = X_{t-k}
  \]

\end{itemize}

\end{frame}

\begin{frame}{Polynômes en $L$}

\begin{itemize}

\item On peut définir des \textbf{polynômes en $L$} :
  \[
  \left(\sum_{i=0}^{p} a_i L^i\right) X_t = \sum_{i=0}^{p} a_i X_{t-i}
  \]\newline

\item On peut aussi définir des \textbf{séries en $L$} (ou $F$). On doit bien sûr se restreindre à des processus stationnaires pour que cela ait un sens.\newline

\item Si $(X_t, t \in \Z)$ est stationnaire au second ordre et $(a_i)_{i \geq 0}$ absolument sommable, alors :
  \[
  Y_t = \sum_{i=0}^{\infty} a_i X_{t-i} = \left(\sum_{i=0}^{\infty} a_i L^i\right) X_t
  \]
  est aussi stationnaire.

\end{itemize}

\end{frame}

\begin{frame}{Propriétés des séries en $L$}

\begin{itemize}

\item \textbf{Addition :}
  \[
  \left(\sum_{i=0}^{\infty} a_i L^i + \sum_{i=0}^{\infty} b_i L^i\right) X_t = \sum_{i=0}^{\infty} (a_i + b_i) L^i X_t
  \]
  où la suite $(a_i + b_i)$ est absolument sommable.\newline

\item \textbf{Multiplication :}
  \[
  \left(\sum_{i=0}^{\infty} a_i L^i\right) \left(\sum_{j=0}^{\infty} b_j L^j\right) X_t = \sum_{k=0}^{\infty} c_k L^k X_t
  \]
  où $c_k = \sum_{i=0}^{k} a_i b_{k-i}$ et la suite $(c_k)$ est absolument sommable.

\end{itemize}

\end{frame}

\begin{frame}{Démonstration de la multiplication (1/2)}

\begin{itemize}

\item Notons $A(L) = \sum_{i=0}^{\infty} a_i L^i$ et $B(L) = \sum_{j=0}^{\infty} b_j L^j$.\newline

\item Calculons $A(L) B(L) X_t$ :
  \begin{align*}
  A(L) B(L) X_t &= A(L) \left( \sum_{j=0}^{\infty} b_j X_{t-j} \right) = \sum_{i=0}^{\infty} a_i \sum_{j=0}^{\infty} b_j X_{t-j-i}
  \end{align*}\newline

\item En posant $k = i + j$ et en réorganisant :
  \[
  A(L) B(L) X_t = \sum_{k=0}^{\infty} \left( \sum_{i=0}^{k} a_i b_{k-i} \right) X_{t-k} = \sum_{k=0}^{\infty} c_k X_{t-k}
  \]
  où $c_k = \sum_{i=0}^{k} a_i b_{k-i}$ est le \textbf{produit de convolution}.

\end{itemize}

\end{frame}

\begin{frame}{Démonstration de la multiplication (2/2)}

\begin{itemize}

\item Il reste à montrer que la suite $(c_k)$ est absolument sommable.\newline

\item On a :
  \begin{align*}
  \sum_{k=0}^{\infty} |c_k| &= \sum_{k=0}^{\infty} \left| \sum_{i=0}^{k} a_i b_{k-i} \right| \leq \sum_{k=0}^{\infty} \sum_{i=0}^{k} |a_i| |b_{k-i}|
  \end{align*}\newline

\item En changeant l'ordre de sommation :
  \[
  \sum_{k=0}^{\infty} |c_k| \leq \left( \sum_{i=0}^{\infty} |a_i| \right) \left( \sum_{j=0}^{\infty} |b_j| \right) < +\infty
  \]
  car $(a_i)$ et $(b_j)$ sont absolument sommables. \hfill $\square$\newline

\item Le produit de deux séries absolument sommables est absolument sommable : l'ensemble des séries en $L$ forme une \textbf{algèbre}.

\end{itemize}

\end{frame}

\begin{frame}{Inversibilité du polynôme $(1 - \lambda L)$}

\begin{itemize}

\item Le polynôme retard $\phi(L) = 1 - \lambda L$ est inversible dès lors que $|\lambda| < 1$.\newline

\item Démonstration : Posons $a_i = \lambda^i$ pour $i \geq 0$. La suite $(a_i)$ est absolument sommable (série géométrique convergente car $|\lambda| < 1$).\newline

\item En multipliant par $\phi(L) = 1 - \lambda L$ :
  \begin{align*}
  (1 - \lambda L) \sum_{i=0}^{\infty} \lambda^i L^i &= \sum_{i=0}^{\infty} \lambda^i L^i - \sum_{i=0}^{\infty} \lambda^{i+1} L^{i+1} \\
  &= 1 + \sum_{i=1}^{\infty} \lambda^i L^i - \sum_{i=1}^{\infty} \lambda^i L^i = 1
  \end{align*}\newline

\item Ainsi $(1 - \lambda L)^{-1} = \sum_{i=0}^{\infty} \lambda^i L^i$. \hfill $\square$

\end{itemize}

\end{frame}

\begin{frame}{Interprétation}

\begin{itemize}

\item Soit $(X_t, t \in \Z)$ un processus stationnaire au second ordre.\newline

\item Alors le processus $(Y_t, t \in \Z)$ défini par :
  \[
  Y_t = \sum_{i=0}^{\infty} \lambda^i X_{t-i}
  \]
  est l'unique processus stationnaire au second ordre solution de l'équation :
  \[
  Z_t - \lambda Z_{t-1} = X_t \quad \Leftrightarrow \quad (1 - \lambda L) Z_t = X_t
  \]\newline

\item Il existe d'autres solutions (une infinité), mais elles ne sont pas stationnaires au second ordre.\newline

\item \textbf{Analogie :} $u_n = \lambda u_{n-1} + b$ a une unique solution constante : $u^* = \frac{b}{1-\lambda}$.

\end{itemize}

\end{frame}

\begin{frame}{Inversion d'un polynôme général}

\begin{itemize}

\item Soit la fonction polynomiale :
  \[
  \phi(z) = 1 + \phi_1 z + \phi_2 z^2 + \cdots + \phi_p z^p
  \]
  dont les racines $z_j = \frac{1}{\lambda_j}$ sont plus grandes que 1 en module.\newline

\item Il existe une série $\psi(z) = \sum_{i=0}^{\infty} \psi_i z^i$ telle que $\phi(z) \cdot \psi(z) = 1$.\newline

\item Le polynôme retard $\phi(L) = 1 + \phi_1 L + \phi_2 L^2 + \cdots + \phi_p L^p$ est inversible et admet pour inverse $\psi(L)$.

\end{itemize}

\end{frame}

\begin{frame}{Méthode d'inversion par identification (1/2)}

\begin{itemize}

\item On cherche $\psi(z) = \sum_{i=0}^{\infty} \psi_i z^i$ tel que $\phi(z) \cdot \psi(z) = 1$.\newline

\item Développons le produit :
  \[
  (1 + \phi_1 z + \phi_2 z^2 + \cdots + \phi_p z^p)(\psi_0 + \psi_1 z + \psi_2 z^2 + \cdots) = 1
  \]\newline

\item En regroupant par puissances de $z$ :
  \begin{align*}
  z^0 &: \quad \psi_0 = 1 \\
  z^1 &: \quad \psi_1 + \phi_1 \psi_0 = 0 \quad \Rightarrow \quad \psi_1 = -\phi_1 \\
  z^2 &: \quad \psi_2 + \phi_1 \psi_1 + \phi_2 \psi_0 = 0 \quad \Rightarrow \quad \psi_2 = -\phi_1 \psi_1 - \phi_2
  \end{align*}

\end{itemize}

\end{frame}

\begin{frame}{Méthode d'inversion par identification (2/2)}

\begin{itemize}

\item \textbf{Formule de récurrence générale :} Pour $k \geq 1$ :
  \[
  \boxed{\psi_k = -\sum_{j=1}^{\min(k,p)} \phi_j \psi_{k-j}}
  \]
  avec la convention $\psi_0 = 1$ et $\phi_j = 0$ pour $j > p$.\newline

\item \textbf{Exemple :} Pour $\phi(L) = 1 - 0.8L$ (AR(1) avec $\phi_1 = -0.8$) :
  \begin{align*}
  \psi_0 &= 1 \\
  \psi_1 &= -(-0.8) \cdot 1 = 0.8 \\
  \psi_2 &= -(-0.8) \cdot 0.8 = 0.64 = 0.8^2 \\
  \psi_k &= 0.8^k
  \end{align*}\newline

\item On retrouve $(1 - 0.8L)^{-1} = \sum_{k=0}^{\infty} 0.8^k L^k$.

\end{itemize}

\end{frame}

\begin{frame}{Cas des racines unitaires : le problème}

\begin{itemize}

\item Que se passe-t-il si le polynôme $\phi(z)$ admet une ou plusieurs racines de module 1 (racines unitaires) ?\newline

\item Si $\phi(z)$ admet une racine $z_0$ avec $|z_0| = 1$, alors le polynôme retard $\phi(L)$ n'est \textbf{pas inversible} au sens usuel : il n'existe pas de série $\psi(L) = \sum_{i=0}^{\infty} \psi_i L^i$ avec $(\psi_i)$ absolument sommable telle que $\phi(L) \psi(L) = 1$.\newline

\item \textbf{Exemple :} Le polynôme $\phi(L) = 1 - L$ a pour racine $z = 1$ (racine unitaire).\newline

\item L'inverse formel serait $\sum_{i=0}^{\infty} L^i$, mais la suite $(1, 1, 1, \ldots)$ n'est pas absolument sommable.

\end{itemize}

\end{frame}

\begin{frame}{Factorisation avec racines unitaires}

\begin{itemize}

\item Supposons que $\phi(z)$ admette $k$ racines unitaires. On peut factoriser :
  \[
  \phi(z) = (1-z)^k \tilde{\phi}(z)
  \]
  où $\tilde{\phi}(z)$ est un polynôme dont toutes les racines sont de module $> 1$.\newline

\item En termes d'opérateur retard :
  \[
  \phi(L) = (1-L)^k \tilde{\phi}(L)
  \]\newline

\item Remarques :
  \begin{itemize}
    \item $(1-L)$ est l'\textbf{opérateur de différenciation} : $(1-L)X_t = X_t - X_{t-1} = \Delta X_t$\newline
    \item $(1-L)^k$ correspond à la différenciation d'ordre $k$ : $\Delta^k X_t$\newline
    \item $\tilde{\phi}(L)$ est inversible car ses racines sont de module $> 1$
  \end{itemize}

\end{itemize}

\end{frame}

\begin{frame}{Inversion partielle}

\begin{itemize}

\item Bien que $\phi(L)$ ne soit pas globalement inversible, on peut écrire :
  \[
  \phi(L) = (1-L)^k \tilde{\phi}(L)
  \]
  et inverser la partie $\tilde{\phi}(L)$ :
  \[
  \tilde{\phi}(L)^{-1} = \tilde{\psi}(L) = \sum_{i=0}^{\infty} \tilde{\psi}_i L^i
  \]
  avec $(\tilde{\psi}_i)$ absolument sommable.\newline

\item Si $\phi(L) Y_t = X_t$, alors :
  \[
  (1-L)^k \tilde{\phi}(L) Y_t = X_t \quad \Rightarrow \quad (1-L)^k Y_t = \tilde{\psi}(L) X_t
  \]\newline

\item On peut exprimer $\Delta^k Y_t$ en fonction de $X_t$, mais pas $Y_t$ directement.

\end{itemize}

\end{frame}

\begin{frame}{Processus intégrés}

\begin{itemize}

\item Un processus $(Y_t)$ est dit \textbf{intégré d'ordre $k$}, noté $Y_t \sim I(k)$, si :
  \begin{itemize}
    \item $Y_t$ n'est pas stationnaire\newline
    \item $\Delta^k Y_t = (1-L)^k Y_t$ est stationnaire
  \end{itemize}

\item \textbf{Exemple :} La marche aléatoire $Y_t = Y_{t-1} + \varepsilon_t$ où $\varepsilon_t \sim BB(0,\sigma^2)$.\newline

\item On a $(1-L)Y_t = \varepsilon_t$, donc $\Delta Y_t = \varepsilon_t$ est stationnaire. Ainsi $Y_t \sim I(1)$.\newline

\item De nombreuses séries économiques (PIB, prix, indices boursiers) sont $I(1)$ : leur niveau n'est pas stationnaire, mais leur taux de croissance l'est.

\end{itemize}

\end{frame}

\begin{frame}{Résolution avec racines unitaires}

\begin{itemize}

\item Si $\phi(L)$ a $k$ racines unitaires, on ne peut pas inverser $\phi(L)$ directement, mais on peut :\newline

\item \textbf{1. Différencier} le processus $k$ fois pour éliminer les racines unitaires :
  \[
  \phi(L) Y_t = X_t \quad \Rightarrow \quad \tilde{\phi}(L) \Delta^k Y_t = X_t
  \]\newline

\item \textbf{2. Inverser} $\tilde{\phi}(L)$ (qui n'a plus de racines unitaires) :
  \[
  \Delta^k Y_t = \tilde{\psi}(L) X_t = \sum_{i=0}^{\infty} \tilde{\psi}_i X_{t-i}
  \]\newline

\item \textbf{3. Intégrer} $k$ fois pour retrouver $Y_t$ (avec $k$ constantes d'intégration).

\end{itemize}

\end{frame}

\begin{frame}{Interprétation des constantes d'intégration}

\begin{itemize}

\item L'intégration introduit $k$ \textbf{constantes d'intégration} qui correspondent aux \textbf{conditions initiales}.\newline

\item \textbf{Exemple pour $k=1$ :} Si $\Delta Y_t = Z_t$ où $Z_t$ est connu, alors :
  \[
  Y_t = Y_0 + \sum_{s=1}^{t} Z_s
  \]
  La constante $Y_0$ est la \textbf{condition initiale}.\newline

\item \textbf{Exemple pour $k=2$ :} Si $\Delta^2 Y_t = Z_t$, alors :
  \[
  Y_t = Y_0 + t \cdot \Delta Y_0 + \sum_{s=1}^{t} (t-s+1) Z_s
  \]\newline

\item Ces conditions initiales déterminent le \textbf{niveau} du processus mais n'affectent pas sa dynamique.

\end{itemize}

\end{frame}

\begin{frame}{Conclusion sur les racines unitaires}

\begin{itemize}

\item Les racines unitaires empêchent l'inversion directe mais permettent une représentation du processus différencié. C'est la base de l'analyse des séries non stationnaires.\newline

\item Points clés :
  \begin{itemize}
    \item Un polynôme avec $k$ racines unitaires ne peut pas être inversé en une série absolument sommable\newline
    \item On factorise : $\phi(L) = (1-L)^k \tilde{\phi}(L)$ où $\tilde{\phi}(L)$ est inversible\newline
    \item Le processus différencié $\Delta^k Y_t$ admet une représentation stationnaire\newline
    \item Les $k$ conditions initiales déterminent le niveau mais pas la dynamique
  \end{itemize}

\end{itemize}

\end{frame}

\begin{frame}{Résumé}

\begin{itemize}

\item Un \textbf{processus stochastique} est une suite de variables aléatoires.\newline

\item La \textbf{stationnarité} (au second ordre) simplifie considérablement l'inférence.\newline

\item La \textbf{fonction d'autocovariance} $\gamma(h)$ caractérise les dépendances.\newline

\item La \textbf{fonction d'autocorrélation} $\rho(h) = \gamma(h)/\gamma(0)$ normalise les dépendances.\newline

\item L'\textbf{autocorrélation partielle} $r(K)$ mesure le lien direct entre $X_t$ et $X_{t-K}$.\newline

\item La \textbf{densité spectrale} $f(\omega)$ donne une vision fréquentielle équivalente.\newline

\item L'\textbf{opérateur retard} $L$ permet une écriture compacte des modèles.

\end{itemize}

\end{frame}

\end{document}
